{\large\textbf{Strongly correlated Fermi gases (10.0 points)}}

Despite twenty orders of magnitude of difference in density, it can be shown
that laser cooled atoms and nuclear matter in neutron star share the same
equation of state characterized by a single numerical parameter called
Bertsch's parameter. In this problem we show how precise measurements on
ultracold vapours allowed for an accurate determination of this parameter using
tabletop experiments.

\begin{center}\includegraphics[width=0.5\linewidth]{figures/IPhO_2025_Q6_fig1.jpg}\end{center}

Fig. 1. Images of a Bose-Einstein condensate of ${}^7\mathrm{Li}$ atoms (top),
immersed in a gas of fermionic ${}^6\mathrm{Li}$ atoms (bottom) both confined in
the same magnetic trap. The condensate (narrow central peak) comprises
$1 \times 10^4$ atoms, and the broad pedestal corresponds to the uncondensed
atoms. The larger axial extension of the fermion cloud ($2.5 \times 10^4$ atoms)
reflects the Fermi pressure resulting from Pauli Principle preventing two
fermions from occupying the same state.

% script E printed as ℰ throughout; \mathcal{E} is the closest symbol in the allowed packages
\textbf{A. Thermodynamics of a non-interacting quantum gas}

Consider a quantum particle of mass $m$ confined in a cubic box of size $L$. We
consider first that the motion of the particle is restricted along the $x$ axis
and we assume that the wave function of the particle can be described in
complex notations by a plane wave
\begin{equation}
\psi(x) = Ae^{ik_x x} + Be^{-ik_x x}. \tag{1}
\end{equation}
Here $k_x$ is positive and we note $\lambda$ the associated wavelength.

\textbf{A.1}\quad Express the kinetic energy $K_x$ of the particle as a function
of $h$, $\lambda$ and $m$. \hfill 0.5pt

Since the particle cannot leave the box, we assume that $\psi(0) = \psi(L) = 0$.

\textbf{A.2}\quad Show that $k_x = k_1 n_x$, where $n_x$ is a strictly positive
integer. \hfill 0.6pt

We assume that the previous result is valid in all three $x$, $y$ and $z$
directions.

\textbf{A.3}\quad Represent the quantum states in the $(k_x, k_y, k_z)$ phase
space. Express with $L$ the volume $(\varDelta k)^3$ occupied by each state.
\hfill 0.7pt

We consider fermionic atoms that, like electrons, obey Pauli's Exclusion
Principle, meaning that we can only put two atoms per quantum state. We assume
that the states are filled one by one with increasing energy.

Consider a number $N \gg 1$ of atoms. We call $\mathcal{E}_F$ (the Fermi energy)
the energy of the last occupied state. We also define the so-called Fermi
momentum $k_F$ by $\mathcal{E}_F = \dfrac{\hbar^2 k_F^2}{2m}$.

\textbf{A.4}\quad Represent the states with energy lower than $\mathcal{E}_F$ in
$\overrightarrow{k}$-space. Deduce the expression of $\mathcal{E}_F$ with $N$.
\hfill 1.4pt

We add a small number of atoms $\mathrm{d}N \ll N$ to the system.

\textbf{A.5}\quad Express the increase in energy $\mathrm{d}\mathcal{E}$ in the
system as a function of $\mathcal{E}_F(N)$. Show that the energy $\mathcal{E}$
is given by
\begin{equation}
\mathcal{E}(N) = \kappa N \mathcal{E}_F(N) \tag{2}
\end{equation}
for some number $\kappa$. \hfill 0.8pt

We note $P$ the pressure of the gas.

\textbf{A.6}\quad Express the energy variation $\mathrm{d}\mathcal{E}$ for a
change in volume $\mathrm{d}V$. Deduce the expression of the pressure as a
function of density $n = N/L^3$. \hfill 0.8pt

We now consider the case of interacting atoms. The interactions are described
by an attractive interatomic potential $V(\overrightarrow{r})$ of typical range
$r_e \simeq a_0$. In the so-called ultra-cold regime, the atomic wavelength is
much larger than $r_e$. As a consequence, the matter waves cannot resolve the
details of the potential and one can show that the effect of the interactions
can be encapsulated in the coupling constant
\begin{equation}
g = \iiint_{\mathbb{R}^3} V\left(\overrightarrow{r}\right) \mathrm{d}^3\overrightarrow{r}. \tag{3}
\end{equation}
Furthermore, we conventionally take $g = \dfrac{4\pi\hbar^2 a}{m}$ , which
defines $a$.

Using a dimensional argument, one can show that the energy of the interacting
cloud can be written as
\begin{equation}
\mathcal{E} = \kappa N \mathcal{E}_F(N)\, f\left(\frac{1}{na^\alpha}\right), \tag{4}
\end{equation}
where $\kappa$ and $\mathcal{E}_F$ were introduced in questions A.4 and A.5, and
$f$ depends only on $\dfrac{1}{na^\alpha}$.

\textbf{A.7}\quad Give the value of $\alpha$. \hfill 0.4pt

The so-called unitary limit corresponds to a regime where $a = \infty$. We
assume that the function $f$ has a finite value in this limit and we define
$\xi = f(0)$ the so-called Bertsch parameter.

\textbf{A.8}\quad Show that, at the unitary limit, the properties of the
interacting cloud are identical to those of a noninteracting system up to a
rescaling of Planck's constant $\hbar \to \xi^\gamma \hbar$. Give the value of
$\gamma$. \hfill 0.6pt

\textbf{B. Thermodynamics of a trapped quantum gas}

We now consider that the atoms are confined by a single-particle optical
potential $U(\overrightarrow{r})$. We assume that the cloud can be locally
considered as homogeneous and that the results of Part A are still applicable
locally. We consider first the case of a non-interacting Fermi gas.

Consider first the case of a one-dimensional potential $U(x)$, with
$U(0) = 0$. We assume that the cloud is homogeeous in the $(y, z)$ direction and
we note $\varSigma$ the area of the cloud in the $(y, z)$ plane. Let's consider
the volume of the cloud comprised between $x$ and $x + \mathrm{d}x$.

\textbf{B.1}\quad Write the total forces exerted on the atoms and show that the
density $n(\overrightarrow{r})$ is given by
\begin{equation}
n(\overrightarrow{r}) = A(\mathcal{E}_F(0) - U(\overrightarrow{r}))^{3/2}, \tag{5}
\end{equation}
where $\mathcal{E}_F(0)$ is the Fermi energy at the trap center. Give the
expression of $A$. \hfill 1.5pt

We assume that this expression holds for a general potential
$U(\overrightarrow{r})$ depending on the three coordinates $(x, y, z)$. If the
the cloud is sufficiently small we can furthermore approximate $U$ by a
harmonic potential close to its minimum. We thus take
\begin{equation}
U(\overrightarrow{r}) = \frac{m}{2} \sum_{i=x,y,z} \omega_i^2 x_i^2. \tag{6}
\end{equation}

\textbf{B.2}\quad Find the equation defining the surface of the cloud. Give the
expression of the radius $R_i$ of the cloud in the direction $i = x, y, z$ as a
function of $\mathcal{E}_F(0)$, $m$ and $\omega_i$. \hfill 0.7pt

Using the previous question, one can prove that after a proper spatial
rescaling, the system can be mapped onto a gas of fermions trapped in an
isotropic harmonic potential of frequency
$\bar{\omega} = (\omega_x \omega_y \omega_z)^{1/3}$.

\textbf{B.3}\quad Calculate the total atom number as a spatial integral by
decomposing the cloud into infinitesimal shells of radius $r$ and thickness
$dr$, and show that
\begin{equation}
\mathcal{E}_F(0) = \hbar\bar{\omega}(3N)^\mu, \tag{7}
\end{equation}
where you will give the value of the exponent $\mu$. \textit{Hint}: we give
\begin{equation}
\int_0^1 r^2(1 - r^2)^{3/2}\mathrm{d}r = \frac{\pi}{32}. \tag{8}
\end{equation}
\hfill 1.1pt

It is possible to change the value of $a$ using an external magnetic field
$\overrightarrow{B}$. In the case of fermionic \textsuperscript{6}Li, the
unitary limit is reached for
$B = \left\|\overrightarrow{B}\right\| = 8.32 \times 10^{-2}\,\mathrm{T}$ and
the cloud behaves as a non-interacting gas at high field.

\begin{center}\includegraphics[width=0.75\linewidth]{figures/IPhO_2025_Q6_fig2.png}\end{center}

Fig. 2. Density profile of a cloud of $7 \times 10^4$ fermionic lithium atoms
for various external magnetic fields (figure from Bourdel \textit{et al.} Phys.
Rev. Lett. \textbf{93}, 050401 (2004)).

\textbf{B.4}\quad Express the ratio $\dfrac{R(a = \infty)}{R(a = 0)}$ with
$\xi$ and deduce from the data of Fig. 2 an estimate of $\xi$. \hfill 0.9pt
